\section{Équation de la chaleur-Semaine 11}

\begin{definition}[Le problème de la chaleur 1D]
Trouver $T : [0, L] \times \mathbb{R}_+ \to \mathbb{R}$ telle que :
\begin{align}
    \rho c_p \frac{\partial T}{\partial t}(x,t) - k \frac{\partial^2 T}{\partial x^2}(x,t) &= f(x,t), \quad 0 < x < L, \ t > 0 \\
    T(0,t) = T(L,t) &= 0 \quad \text{(Conditions aux limites)} \label{eq:cl} \\
    T(x,0) &= w(x), \quad 0 < x < L \quad \text{(Condition initiale)} \label{eq:ci}
\end{align}
\end{definition}

\begin{center}
\begin{tikzpicture}[xscale=1.5]
    % Barre
    \draw[thick] (0,0) -- (6,0);
    \fill[gray!20] (0,-0.1) rectangle (6,0.1);
    \draw[thick] (0,-0.2) -- (0,0.2) node[above] {$0$};
    \draw[thick] (6,-0.2) -- (6,0.2) node[above] {$L$};
    \node[below] at (3, -0.2) {Axe $x$};
    
    % Source de chaleur
    \node[algofr] at (3, 0.8) {Source de chaleur $f(x,t)$};
    \foreach \x in {1.5, 3, 4.5} {
        \draw[->, algofr, thick] (\x, 0.6) -- (\x, 0.15);
    }
    
    % Conditions
    \node[deffr, align=center] at (0, -0.6) {$T=0$};
    \node[deffr, align=center] at (6, -0.6) {$T=0$};
\end{tikzpicture}
\end{center}

\begin{explication}[Données du problème]
\begin{itemize}
    \item $T(x,t)$ : Température au point $x$ et au temps $t$.
    \item $w(x)$ : Température initiale connue.
    \item $\rho > 0$ : Densité.
    \item $c_p > 0$ : Chaleur spécifique.
    \item $k > 0$ : Conductivité thermique.
    \item $f$ : Source de chaleur (Puissance par unité de longueur).
\end{itemize}
On cherche la température $T(x,t)$ pour $x \in [0,L]$ et $t > 0$.
\end{explication}

\begin{remarque}[Principe de conservation]
L'équation de la chaleur traduit le principe de conservation de l'énergie thermique. Si on intègre l'équation par rapport à $x$ entre $0$ et $L$, on obtient :
\[
    \underbrace{\int_0^L \rho c_p \frac{\partial T}{\partial t}(x,t) dx}_{\text{Variation de l'énergie thermique}} - \underbrace{\int_0^L k \frac{\partial^2 T}{\partial x^2}(x,t) dx}_{\text{Flux de chaleur}} = \underbrace{\int_0^L f(x,t) dx}_{\text{Puissance thermique}}
\]
Ce qui se réécrit :
\[
    \frac{d}{dt} \int_0^L \rho c_p T(x,t) dx - \left[ k \frac{\partial T}{\partial x}(x,t) \right]_{x=0}^L = \int_0^L f(x,t) dx
\]
\end{remarque}

\subsection{Problème modèle}

Quitte à diviser toute l'équation par $\rho c_p$, on peut se ramener à considérer l'équation normalisée suivante (en redéfinissant $k \leftarrow k/(\rho c_p)$ et $f \leftarrow f/(\rho c_p)$) :

\begin{equation}
\begin{cases}
    \frac{\partial u}{\partial t}(x,t) - k \frac{\partial^2 u}{\partial x^2}(x,t) = f(x,t), & 0 < x < L, \ t > 0 \\
    u(0,t) = u(L,t) = 0, & t > 0 \\
    u(x,0) = w(x), & 0 < x < L
\end{cases}
\end{equation}

\begin{exemple}[Solution exacte analytique]
Si $f(x,t) = 0$ et la condition initiale est $w(x) = \sin\left(\frac{\pi x}{L}\right)$, alors la solution exacte est donnée par :
\[
    u(x,t) = \sin\left(\frac{\pi x}{L}\right) \exp\left(-\frac{k \pi^2}{L^2} t\right)
\]
\end{exemple}

\begin{proposition}[Principe du maximum continu]
Supposons $f(x,t) = 0$ et soit $u$ la solution de l'équation homogène. Alors :
\begin{enumerate}
    \item Si $w(x) \ge 0$, alors $u(x,t) \ge 0$ pour tout $t \ge 0$.
    \item $\max_{x \in [0,L]} u(x,t) \le \max_{x \in [0,L]} w(x)$.
\end{enumerate}
\end{proposition}

\begin{center}
\begin{tikzpicture}
    \begin{axis}[
        axis lines = middle,
        xlabel = {$x$},
        ylabel = {$u(x,t)$},
        xmin=-0.1, xmax=1.1,
        ymin=-0.1, ymax=1.5,
        xtick={1}, xticklabels={$L$},
        ytick=\empty,
        width=10cm, height=6cm,
        legend pos=north east,
        legend style={draw=none, fill=none}
    ]
    % Condition initiale w(x)
    \addplot[domain=0:1, samples=100, thick, black] {4*x*(1-x) + 0.2*sin(180*x)};
    \addlegendentry{$u(x,0) = w(x)$}
    
    % Solution u(x,t)
    \addplot[domain=0:1, samples=100, thick, deffr] {2.5*x*(1-x) + 0.1*sin(180*x)};
    \addlegendentry{$u(x,t)$ pour $t > 0$}
    
    \node[below left] at (axis cs:0,0) {$0$};
    \node[above, text=black] at (axis cs:0.5, 1.25) {$u(x,0) \ge u(x,t)$};
    \end{axis}
\end{tikzpicture}
\end{center}

\subsection{Approximation numérique}

On cherche à approximer la solution $u$ de l'EDP pour $(x,t) \in [0,L] \times [0,T]$.
Pour cela, on discrétise le domaine spatio-temporel.

\textbf{Discrétisation en espace :}
$0 = x_0 < x_1 < \dots < x_{N+1} = L$, avec un pas $h = \frac{L}{N+1}$ tel que $x_i = i \cdot h$.

\textbf{Discrétisation en temps :}
$0 = t_0 < t_1 < \dots < t_M = T$, avec un pas $\tau = \frac{T}{M}$ tel que $t_n = n \cdot \tau$.

\begin{center}
\begin{tikzpicture}[scale=1.5]
    \draw[->, thick] (-0.2,0) -- (5.5,0) node[right] {$x$};
    \draw[->, thick] (0,-0.2) -- (0,3.5) node[above] {$t$};
    
    % Grid lines
    \foreach \x in {1,2,3,4} \draw[gray!50, very thin] (\x,0) -- (\x,3);
    \foreach \y in {0.6,1.2,1.8,2.4,3.0} \draw[gray!50, very thin] (0,\y) -- (5,\y);
    
    % Nodes X
    \foreach \x in {0,1,2,3,4,5} \draw[thick] (\x,0.05) -- (\x,-0.05);
    \node[below] at (0,0) {$0$};
    \node[below] at (1,0) {$x_1$};
    \node[below] at (2,0) {$x_2$};
    \node[below] at (3,0) {$\dots$};
    \node[below] at (4,0) {$x_N$};
    \node[below, align=center] at (5,0) {$L$\\$=x_{N+1}$};
    
    % Nodes T
    \foreach \y in {0.6,1.2,1.8,2.4,3.0} \draw[thick] (0.05,\y) -- (-0.05,\y);
    \node[left] at (0,0.6) {$t_1$};
    \node[left] at (0,1.2) {$t_2$};
    \node[left] at (0,1.8) {$\dots$};
    \node[left] at (0,3.0) {$T = t_M$};

    % Indication condition initiale
    \draw[<-, thmfr, thick] (4.5,0.1) -- (5, 0.8) node[right, align=left] {$u(x,0)$ connu\\(Condition initiale)};
\end{tikzpicture}
\end{center}

On cherche à trouver des approximations $u_i^n \approx u(x_i, t_n)$ avec l'indice $n$ pour le temps et l'indice $i$ pour l'espace, pour $i=1,\dots,N$ et $n=1,\dots,M$.

\textbf{Initialisation :}
On connaît la condition initiale $u(x_i, t_0) = u(x_i, 0) = w(x_i)$ pour $i=1,\dots,N$.
On pose donc : $u_i^0 = w(x_i)$.

\textbf{Question :} Comment calculer $u_i^{n+1}$ si on connaît les valeurs au pas de temps précédent $u_i^n$ ?

\subsubsection{Schéma d'Euler progressif (explicite)}

On fixe $n$ et on réécrit l'équation de la chaleur évaluée au point $(x_i, t_n)$ :
\[
    \frac{\partial u}{\partial t}(x_i,t_n) - k \frac{\partial^2 u}{\partial x^2}(x_i,t_n) = f(x_i,t_n)
\]
On approche les dérivées par différences finies (progressives en temps, centrées en espace) :
\[
    \underbrace{\frac{u(x_i, t_{n+1}) - u(x_i, t_n)}{\tau}}_{\text{Différences finies progressives}} - k \underbrace{\frac{u(x_{i+1}, t_n) - 2u(x_i, t_n) + u(x_{i-1}, t_n)}{h^2}}_{\text{Différences finies centrées}} = f(x_i, t_n) + \mathcal{O}(\tau + h^2)
\]
En remplaçant la solution exacte $u(x_i, t_n)$ par notre approximation numérique $u_i^n$ et en tronquant l'erreur $\mathcal{O}(\tau + h^2)$, on obtient le schéma explicite suivant :

\begin{algorithme}[Méthode à un pas en temps (Euler Explicite)]
\begin{equation}
\begin{cases}
    \frac{u_i^{n+1} - u_i^n}{\tau} - k \frac{u_{i+1}^n - 2u_i^n + u_{i-1}^n}{h^2} = f(x_i, t_n) \\
    u_0^n = u_{N+1}^n = 0 \quad \text{(Conditions aux limites)}
\end{cases}
\end{equation}
En isolant $u_i^{n+1}$, la relation de récurrence devient :
\[
    u_i^{n+1} = \left( 1 - \frac{2k\tau}{h^2} \right) u_i^n + \frac{k\tau}{h^2} \left( u_{i+1}^n + u_{i-1}^n \right) + \tau f(x_i, t_n)
\]
\end{algorithme}

\begin{center}
\begin{tikzpicture}[scale=1.5]
    \draw[->, thick] (0,0) -- (4,0) node[right] {$x$};
    \draw[->, thick] (1,-0.5) -- (1,2.5) node[above] {$t$};

    \draw[dashed, gray!50] (1.5,0.5) -- (3.5,0.5);
    \draw[dashed, gray!50] (1.5,1.5) -- (3.5,1.5);
    \draw[dashed, gray!50] (2,0) -- (2,2);
    \draw[dashed, gray!50] (2.5,0) -- (2.5,2);
    \draw[dashed, gray!50] (3,0) -- (3,2);

    % Nœuds connus (temps n)
    \filldraw[deffr] (2,0.5) circle (2pt) node[below left, text=black] {$u_{i-1}^n$};
    \filldraw[deffr] (2.5,0.5) circle (2pt) node[below right, text=black] {$u_i^n$};
    \filldraw[deffr] (3,0.5) circle (2pt) node[below right, text=black] {$u_{i+1}^n$};

    % Nœud à calculer (temps n+1)
    \filldraw[remfr] (2.5,1.5) circle (2pt) node[above right, text=black] {$u_i^{n+1}$};

    % Flèches de dépendance
    \draw[->, thick, shorten >=4pt] (2,0.5) -- (2.5,1.5);
    \draw[->, thick, shorten >=4pt] (2.5,0.5) -- (2.5,1.5);
    \draw[->, thick, shorten >=4pt] (3,0.5) -- (2.5,1.5);

    \node[left] at (1,0.5) {$t_n$};
    \node[left] at (1,1.5) {$t_{n+1}$};
    \node[below] at (2,0) {$x_{i-1}$};
    \node[below] at (2.5,0) {$x_i$};
    \node[below] at (3,0) {$x_{i+1}$};
\end{tikzpicture}
\end{center}

Sous forme matricielle, en posant $U^n = (u_1^n, \dots, u_N^n)^T$ et $F(t_n) = (f(x_1, t_n), \dots, f(x_N, t_n))^T$, le schéma s'écrit :
\[
    U^{n+1} = (I - \tau A) U^n + \tau F(t_n)
\]
où $A$ est la matrice tridiagonale :
\[
A = \frac{k}{h^2}
\begin{pmatrix}
2 & -1 & & 0 \\
-1 & \ddots & \ddots & \\
& \ddots & 2 & -1 \\
0 & & -1 & 2
\end{pmatrix} \in \mathbb{R}^{N \times N}
\]

\subsubsection{Schéma d'Euler rétrograde (implicite)}

On fixe $n$ et on écrit cette fois l'équation au point $(x_i, t_{n+1})$ :
\[
    \frac{\partial u}{\partial t}(x_i, t_{n+1}) - k \frac{\partial^2 u}{\partial x^2}(x_i, t_{n+1}) = f(x_i, t_{n+1})
\]
En appliquant des différences finies rétrogrades pour le temps et centrées pour l'espace, nous obtenons :
\[
    \frac{u(x_i, t_{n+1}) - u(x_i, t_n)}{\tau} - k \frac{u(x_{i+1}, t_{n+1}) - 2u(x_i, t_{n+1}) + u(x_{i-1}, t_{n+1})}{h^2} = f(x_i, t_{n+1}) + \mathcal{O}(\tau + h^2)
\]

\begin{algorithme}[Méthode d'Euler Implicite]
On considère donc le schéma numérique suivant :
\begin{equation}
\begin{cases}
    \frac{u_i^{n+1} - u_i^n}{\tau} - k \frac{u_{i+1}^{n+1} - 2u_i^{n+1} + u_{i-1}^{n+1}}{h^2} = f(x_i, t_{n+1}) \\
    u_0^{n+1} = u_{N+1}^{n+1} = 0
\end{cases}
\end{equation}
En réarrangeant les termes pour isoler les inconnues au temps $t_{n+1}$ :
\[
    \left( 1 + \frac{2k\tau}{h^2} \right) u_i^{n+1} - \frac{k\tau}{h^2} \left( u_{i+1}^{n+1} + u_{i-1}^{n+1} \right) = u_i^n + \tau f(x_i, t_{n+1})
\]
\end{algorithme}

Sous forme matricielle, le système devient :
\[
    (I + \tau A) U^{n+1} = U^n + \tau F(t_{n+1})
\]
À chaque itération, il faut résoudre un système linéaire pour trouver le vecteur $U^{n+1}$ !

\subsection{Stabilité}

Dans ce qui suit, on suppose qu'il n'y a pas de terme source ($f=0$).

\begin{definition}[Critères de Stabilité]
Un schéma numérique qui approxime la solution de l'équation de la chaleur est stable si :
\begin{enumerate}
    \item \textbf{Principe du maximum discret (positivité) :} $\forall i=1,\dots,N, \ u_i^n \ge 0 \implies \forall i=1,\dots,N, \ u_i^{n+1} \ge 0$.
    \item \textbf{Non-croissance de la norme infini :} $\max_{i=1,\dots,N} |u_i^{n+1}| \le \max_{i=1,\dots,N} |u_i^n|$.
\end{enumerate}
On aimerait en effet que l'approximation numérique reproduise le comportement physique de la solution exacte continue.
\end{definition}

\begin{theoreme}[Stabilité des schémas d'Euler]
\begin{itemize}
    \item On peut montrer que le \textbf{schéma d'Euler rétrograde (implicite) est inconditionnellement stable} (stable quels que soient les choix de $h$ et $\tau$).
    \item Le \textbf{schéma d'Euler progressif (explicite) est stable sous condition CFL} : il est stable si et seulement si 
    \[
        \tau \le \frac{h^2}{2k}
    \]
\end{itemize}
\end{theoreme}

\subsection{Convergence}

On s'intéresse à l'erreur d'approximation à l'instant final $T = t_M$ :
\[
    e := \max_{i=1,\dots,N} |u(x_i, T) - u_i^M|
\]
Pour les deux méthodes (Euler progressif / rétrograde), on a utilisé :
\begin{itemize}
    \item 1 formule d'ordre 1 en temps : $\mathcal{O}(\tau)$.
    \item 1 formule d'ordre 2 en espace : $\mathcal{O}(h^2)$.
\end{itemize}

\begin{proposition}[Ordre des schémas]
On peut montrer que les deux schémas vus dans ce cours sont d'ordre $\mathcal{O}(\tau + h^2)$.

Autrement dit : si on divise le pas d'espace $h$ par $2$ et le pas de temps $\tau$ par $4$, l'erreur globale est divisée par $4$.
\end{proposition}